\documentclass{article}
\usepackage[left=2cm, right=2cm, top=2cm, bottom=2cm]{geometry}
\usepackage{graphicx}
\usepackage{amsmath}
\usepackage{amssymb}
\usepackage{amsthm}
\usepackage{fancyhdr}
\usepackage{verbatim}
\usepackage{listings}
\usepackage{xcolor}
\usepackage{pdfpages}
\usepackage{cancel}
\usepackage{physics}
\usepackage{esint}

\lstset{
    language=C++,
    basicstyle=\ttfamily\footnotesize,
    keywordstyle=\color{blue},
    commentstyle=\color{green},
    stringstyle=\color{red},
    numbers=left,
    numberstyle=\tiny\color{gray},
    frame=single,
    breaklines=true,
    captionpos=b,
}


\title{MTH 553: Lecture \# 29}
\author{Cliff Sun}

\newtheorem{theorem}{Theorem}[section]
\newtheorem{lemma}[theorem]{Lemma}
\newtheorem{definition}[theorem]{Definition}
\newtheorem{conjecture}[theorem]{Conjecture}
\newtheorem{proposition}[theorem]{Proposition}
\newtheorem{corollary}[theorem]{Corollary}
\newtheorem{one minute paper}[theorem]{One Minute Paper}

\pagestyle{fancy}
\lhead{\textbf{\thepage}\ \ \nouppercase{\rightmark}}
\chead{MTH 553: Lecture \# 29}
\rhead{Cliff Sun}

\begin{document}

\maketitle

\section*{Lecture Span}
\begin{enumerate}
    \item Distributions
\end{enumerate}

Recall last time, we defined what a derivative of a distribution was. Examples, let $$H_y(x) = \begin{cases}
    0, & x \leq y \\
    1, & x > y
\end{cases}$$

We proved that 

\begin{equation}
    H'(y) = \delta(y)
\end{equation}

Another example, let $$u_y(x) = \frac{1}{2}|x-y|$$

Then, 

\begin{equation}
    u_y'' = \delta_y
\end{equation}

The first derivative gives the step function, then the second derivative gives the delta function. 

\begin{proof}
    \begin{equation}
        u'_y(x) = \begin{cases}
            -1/2, & x < y \\
            1/2, & x > y
        \end{cases} = H_y(x) - 1/2
    \end{equation}
    Then 
    \begin{equation}
        u''_y  = \delta_y
    \end{equation}
\end{proof}

\section{Fundamental Solutions}

Let $L$ be a linear differential operator, then let some region $\Omega$. Then consider a point $y$ with a $\delta_y$. Then solve 

\begin{equation}
    Lu = \delta_y
\end{equation}

For each fixed $y \in \Omega$. Suppose we have $G(x,y)$ which is locally integrable with respect to $x$ which also satisfies 

\begin{equation}
    LG(x,y) = \delta_y(x)
\end{equation}

This is simply the greens function. I.e. 

\begin{align*}
    \int_{\Omega}G(x,y)L'v(x) dx &= \int_{\Omega}\delta_y v(y)dy \\
    &= v(y), \quad \forall v \in C^{\infty}_{0}(\Omega)
\end{align*}

Then, the function 

\begin{equation}
    w(x) = \int_\Omega G(x,y)f(y)dy
\end{equation}

Solves 

\begin{equation}
    Lw = f
\end{equation}

Here, $G(x,y)$ is the fundamental solution to $L$. Note, $w$ is analagous to the "sum" of responses to the impulses. 

\begin{proof}
    We prove that $Lw = f$ weakly. Let $w \in L^1$. (also, we can swap integration variables) 
    \begin{align}
        \int_{\Omega} w(x)L'v(x) dx &= \int_{\Omega}\int_{\Omega}G(x,y)Lv'(x)dx f(y)dy \\
        &= \int_{\Omega}v(y)f(y)dy, \quad \text{because } LG = \delta_y
    \end{align}
    Therefore, $Lw = f$ weakly. 
\end{proof}

For example, $L = d_x^2$ on $\Omega$. THen $$G(x,y) = -\frac{1}{2}|x-y|$$

is a greens function. This is because 

\begin{equation}
    d_x^2 G(x,y) = \delta_y
\end{equation}

On a bounded domain $\Omega$, then we get a Greens function from Dirichlet eigenfunctions. 

\begin{equation}
    G(x,y) = \sum_k \frac{1}{\lambda_k}\phi_k(x)\phi_k(y)
\end{equation}

Note that $G$ is symmetric under exchange of variables.

\begin{proof}
    Proof of eigenfunction expansion of the greens function. For a test function $v$, 
    \begin{align*}
        \int_{\Omega}G(x,y)(-\triangle v)(x)dx &= \sum_k \frac{1}{\lambda_k}\int_{\Omega}\phi_k(x)(-\triangle v)dx\phi_k(y)dy \\
        &= \sum_k \frac{1}{\lambda_k}\int_{\Omega}(-\triangle \phi_k(x))v(x)dx\phi_k(y)dy \\
        &= \sum_k \frac{1}{\lambda_k}\int_{\Omega}(\lambda_k \phi_k(x))v(x)dx\phi_k(y)dy \\
        &= \sum_k \int_{\Omega}(\phi_k(x))v(x)dx\phi_k(y)dy \\
        &= \sum_k \langle \phi_k, v\rangle_{L^2}\phi_k(y)dy \\
        &= v(y)
    \end{align*}
    This proves that 
    \begin{equation}
        -\triangle G = \delta_y
    \end{equation}
    as distributions applied to the test function. Consequence, 
    \begin{equation}
        w(x) = \int_{\Omega}G(x,y)f(y)dy = \sum_k \frac{a_k}{\lambda_k}\phi_k(x)
    \end{equation}
    Where each
    \begin{equation}
        a_k = \int_{\Omega}\phi_k(y)f(y)dy = \langle \phi_k, f \rangle_{L^2}
    \end{equation}
\end{proof}

\section*{Fundamental Solution on $\mathbb{R}^n$ and potentials}

Motivation: we want $K(x)$ solving 

\begin{equation}
    -\triangle K = \delta_0
\end{equation}

Call $K$ the fundamental solution for Poisson's equation on $\mathbb{R}^n$ with pole at the origin. 
Physically, $K$ is the gravitational potential with a point mass at the origin. Because the Laplacian is rotationally invariant, $K$ should also be rotationally invariant. 

Guess $K = \psi(r)$:

\begin{align}
    \psi''(r) + \frac{n-1}{r}\psi'(r) &= 0
\end{align}

For all points in space except for $x \neq 0$. Here, 

\begin{align*}
    (r^{n-1}\psi'(r))' &= 0 \\
    \psi'(r) &= \frac{c}{r^{n-1}}
\end{align*}

Therefore, 

\begin{equation}
    \psi(r) = \begin{cases}
    c_1 + c_2 r, & n=1 \\
    c_1 + c_2\log(1/r), & n=2 \\
    c_1 + c_2\frac{1}{r^{n-2}}, & n \geq 3
    \end{cases}
\end{equation}

Take $c_1 = 0$, since it doesn't affect the Laplacian. Choose $c_2$ such that 

\begin{equation}
    -\triangle K = \delta_0
\end{equation}



\end{document}